\section{Introduction}
\label{sec:introduction}

Scientific inference requires both executing a statistical procedure and
deciding whether its assumptions are supported.  A numerically correct answer
can remain scientifically invalid when dependence, instability, or
distributional shape makes the selected procedure inadmissible; abstention may
then be preferable to a plausible estimate.  Although language models
increasingly assist scientific and statistical analysis
\citep{GuEtAl2024BLADE,LiuEtAl2024QRData}, evaluations more often test
calculation than \emph{assumption-aware method selection}: connecting observed
diagnostic evidence to a justified procedure.

Structured workflows are natural candidates for this decision, but their
comparisons are confounded.  Moving from a concise prompt to an ``agentic''
system may simultaneously enrich instructions, add calls, preserve
intermediate state, partition context, and introduce role-specific prompts
\citep{YaoEtAl2023ReAct,HongEtAl2024MetaGPT}.  An aggregate gain therefore
does not identify whether instruction quality, staged execution, or role
specialization caused the difference.

We use \ERI{} as a controlled study setting for Sharpe-ratio method selection.
ERI maps data to diagnostic evidence, then to method eligibility, and finally
to a primary procedure or explicit abstention.  The 18 matched cells combine
two time-series families with three information regimes.  Persistent GARCH
returns require procedures robust to serial dependence, whereas a structural
mean break makes the registered full-sample target undefined and requires
abstention.  This contrast makes validity failures observable even when a
workflow produces fluent reasoning or a plausible numerical answer.

Sharpe-ratio inference provides a compact but nontrivial test bed because the
same point estimate admits materially different uncertainty procedures.
Serial dependence and conditional heteroskedasticity can invalidate IID
procedures; non-Gaussian shape changes variance approximations without
removing their IID basis; and temporal instability can invalidate the
full-sample target itself.  A system must therefore distinguish a weakened
method from an ineligible one and recognize when no numerical procedure can
repair an ill-defined target.  The task rewards a defensible evidential chain,
not merely selection of a sophisticated estimator.

Four frozen
conditions progress from concise two-call prompting (A), to enriched unified
two-call instructions (B), to a unified four-call implementation with the same
substantive instructions (C), and finally to a specialized four-role workflow
(D).  This sequence gives the cleanest available comparison of instruction
enrichment (\(B-A\)), a qualified comparison of call organization
(\(C-B\)), and a residual workflow comparison (\(D-C\)).

Three findings organize the paper.  First, richer instructions produce the
most stable improvement.  Second, C has the strongest aggregate profile, but
its advantage over B is concentrated in one severe error and is not robust to
trimming or scenario family.  Third, D performs worse than C under the tested
bundles, which does not isolate persona identity or establish that multi-agent
systems are generally inferior.  As a secondary boundary result, a
domain-specific deterministic policy performs strongly when the benchmark's
evidence and admissibility rules are encoded directly.

Our contribution is a component-sensitive evaluation of instruction richness,
unified multi-call execution, and specialized-role workflow design on a
validity-sensitive benchmark.  We retain failures in the denominator,
measure invalid selection, harmful non-abstention, claim validity, and routing
regret, and report concentration and sensitivity analyses that prevent the
best aggregate result from being overinterpreted.  Typed evidence, eligibility
rules, abstention, audit trails, and deterministic routing support this study;
they are not presented as independent architectural novelties.

This design addresses two common interpretive problems.  First, matching
cells and nominal condition-level output capacity prevents a four-call
workflow from being credited simply for receiving a larger output budget.
Second, complete denominators and fail-closed accounting prevent provider or
schema failures from disappearing from condition averages.  Cell-level
contrasts, scenario-family summaries, completed-pair analyses, and trimming
then reveal whether an aggregate gain is diffuse or concentrated in one
consequential decision.  More broadly, the study shows that aggregate workflow
rankings can be misleading unless instruction content is matched, terminal
failures remain in the denominator, and paired cell-level analyses reveal
whether gains are diffuse or concentrated in rare high-cost errors.

This framing differs from asking whether one agent category is globally
better.  The conditions form a sequence of increasingly structured workflow
bundles, not a complete factorial design.  We identify the cleanest contrast,
state the extra factors changed by the remaining contrasts, and limit
conclusions to the frozen pilot.  The analysis therefore centers on
component-sensitive evaluation and failure modes rather than proposing a new
agent architecture or execution harness.

The remainder of the paper positions this evaluation problem, defines the
evidence-routing task and benchmark, describes the frozen A--D sequence,
reports canonical and sensitivity results, and closes with design
implications and limitations.  The supplement provides schemas, diagnostic
definitions, cell registries, failure reconstruction, and reproducibility
records.
